% ============================================================
% frq-template.tex  (parameterized)
% ============================================================

% ----------------------------
% Helpers for FRQ 23 MC table
% ----------------------------
\providecommand{\mca}[1]{\number\numexpr \MCa*(#1)\relax}
\providecommand{\mcb}[1]{\number\numexpr \MCb*(#1)\relax}

% ============================================================
% Free Response block (styled like MC questions)
% Assumes your main template defines:
%   - \newlist{questions}{enumerate}{1}
%   - \newcommand{\question}{...}
% ============================================================

\emph{Answer the problems on the back of your bubble sheet.}

\begin{questions}
    \setcounter{questionsi}{21} % so the next \item in \begin{questions} prints 22.
    \question{Consider the market for widgets, with demand and supply given by:
        \[
            Q_D = \QDintercept - \QDslope P
            \qquad
            Q_S = \QSintercept + \QSslope P
        \]
        Answer the following:
        \begin{enumerate}[label=\alph*), itemsep=0.35em, topsep=0.2em]
            \item Sketch the supply and demand diagram. Be sure to label the \emph{price intercepts}
                  (i.e., the intercepts on the vertical axis, where \(Q=0\)) of both curves.
            \item Solve for the equilibrium price, \( P^* \). Show your work.
            \item Suppose the government imposes a per-unit tax of \( \Tax \) on sellers.
                  What is the new quantity with the tax, \( Q_t \)?
            \item With the tax in place, what price do buyers pay, \(P_b\)?
            \item Does the tax burden fall more on buyers or sellers, or is it split evenly?
                  (Compare the change in \(P_b\) to the change in the price sellers keep, \(P_s\).)
                  Briefly explain.
        \end{enumerate}
    }

    \question{Two firms, A and B, produce the same good. Their marginal costs are shown in the table below.
        \begin{center}
            \begin{tabular}{c|cc}
                \toprule
                $Q$ & $MC_A$  & $MC_B$  \\
                \midrule
                0   & \mca{0} & \mcb{0} \\
                1   & \mca{1} & \mcb{1} \\
                2   & \mca{2} & \mcb{2} \\
                3   & \mca{3} & \mcb{3} \\
                4   & \mca{4} & \mcb{4} \\
                5   & \mca{5} & \mcb{5} \\
                \bottomrule
            \end{tabular}
        \end{center}

        \begin{enumerate}[label=\alph*), itemsep=0.35em, topsep=0.2em]
            \item First suppose the market price is \( P=\PriceFRQ \). How many units does each firm supply?
                  (Assume whole-number quantities. As usual, each firm produces when \(MC \le P\) and stops when the next unit would have \(MC>P\).)
            \item Instead suppose the total quantity to be produced is fixed at \( Q_A + Q_B = \QTotal \).
                  If production is efficient (total cost minimized), what are \( Q_A \) and \( Q_B \)?
                  Assume quantities must be whole numbers. Explain briefly.
        \end{enumerate}
    }

\end{questions}